\section{Distinguish RSA primitives from RSA schemes}
The calculation \(m^e\bmod n\) is an RSA arithmetic primitive when its representative
and key satisfy the relevant bounds. It does not define how an arbitrary message is
encoded, how randomness is used, or how errors are reported. The toy calculation in
this chapter deliberately exposes the integer operation so its carries and exponent
bits can be checked.

RFC 8017 specifies RSA encryption and signature schemes around these primitives.
RSAES-OAEP encodes a message for encryption using a hash function and mask generation;
RSASSA-PSS specifies an encoding for signatures with its own parameters and verification
rules. Their encodings belong to their respective schemes. Adding an informal padding
byte to the toy integer would not implement either scheme. See the RFC's sections 5,
7.1 and 8.1 for the precise boundaries.

The current routines also branch on operand and exponent bits. They have not been
designed or evaluated for resistance to timing or other side channels. The fixed
128-bit teaching range, toy primes and visible intermediate states support arithmetic
experiments. They do not justify using this code with private keys or real messages.
A production application needs a reviewed cryptographic implementation of its chosen
scheme and key-handling requirements.

Before leaving the chapter, explain where each width enters: a 32-bit limb, a 64-bit
partial-product accumulator, a 129-bit sum represented by four words plus carry, a
256-bit full product, and a 128-bit exponent scan. For each, identify which bits would
be lost by substituting a word operation. That explanation is a useful test of whether
an optimization preserves the calculation.
